% Figaro Paper B: From Firms to Transaction-Scoped Institutions
% A Coasean Re-Examination Under Costless Bilateral Enforcement
% Preprint, April 2026

\documentclass[11pt,a4paper]{article}

% ── Packages ────────────────────────────────────────────────────────
\usepackage[utf8]{inputenc}
\usepackage[T1]{fontenc}
\usepackage{amsmath,amssymb,amsthm}
\usepackage{mathtools}
\usepackage{booktabs}
\usepackage{graphicx}
\usepackage{hyperref}
\usepackage{cleveref}
\usepackage{enumitem}
\usepackage{xcolor}
\usepackage[margin=1in]{geometry}
\usepackage{natbib}
\bibliographystyle{plainnat}

% ── Theorem environments ────────────────────────────────────────────
\newtheorem{theorem}{Theorem}[section]
\newtheorem{proposition}[theorem]{Proposition}
\newtheorem{corollary}[theorem]{Corollary}
\theoremstyle{definition}
\newtheorem{definition}[theorem]{Definition}
\theoremstyle{remark}
\newtheorem{remark}[theorem]{Remark}

% ── Macros ──────────────────────────────────────────────────────────
\newcommand{\figaro}{\textsc{Figaro}}
\newcommand{\pay}{P}
\newcommand{\cumval}{G}

% ────────────────────────────────────────────────────────────────────

\title{%
  From Firms to Transaction-Scoped Institutions:\\
  A Coasean Re-Examination Under Costless Bilateral Enforcement%
}

\author{
  Alessandro Daliana\thanks{Corresponding author. \texttt{adaliana@figaro.org}}
}

\date{April 2026}

\begin{document}

\maketitle

% ════════════════════════════════════════════════════════════════════
\begin{abstract}
\citet{coase1937} argued that firms exist because the cost of
discovering prices, negotiating terms, and enforcing contracts in open
markets sometimes exceeds the cost of hierarchical management.
\citet{williamson1985} refined this into Transaction Cost Economics
(TCE), identifying asset specificity, uncertainty, and frequency as
drivers of vertical integration. Both arguments rest on an implicit
assumption: that contract enforcement between strangers is costly,
slow, and jurisdiction-dependent. This paper examines what happens
when that assumption no longer holds.

We take as given the existence of a settlement primitive---developed
formally in the companion mechanism paper, with a
verified Solidity implementation in the companion implementation
paper---that prices
the cost of trust at exactly $2\times$ the transaction value in
temporarily locked capital, with no recurring extraction, no
arbitrator, and no jurisdictional dependency. Treating this primitive
as exogenous, we ask three questions. First, what does the existence
of such a primitive imply for the Coasean threshold that defines the
boundary of the firm? Second, which functions of the firm dissolve
under it, and which persist in revised form? Third, why is the
implication structurally invisible---and what does that invisibility
itself say about the architecture of contemporary coordination
infrastructure?

Our argument is that the firm bundles several functions historically
inseparable in practice but logically distinct: coordination of
production, capital aggregation, organizational knowledge,
risk pooling, and legal personhood. A costless bilateral enforcement
primitive dissolves the first while leaving the others largely
intact. The result is not the abolition of durable institutions but a
structural shift: the \emph{coordination firm}, whose value is
internal management of a bounded labor pool, loses its economic
rationale where the primitive is applicable; the \emph{capital firm},
whose value is durable asset ownership and risk-bearing, persists in
revised wrappers. We close with a Gramscian observation about why
this shift is harder to see than to specify.
\end{abstract}

\medskip
\noindent\textbf{Keywords:}
transaction cost economics, theory of the firm, smart contracts,
coordination economics, post-firm organization, cultural hegemony

% ════════════════════════════════════════════════════════════════════
\section{Introduction}
\label{sec:intro}

The boundary of the firm has been the central object of theoretical
attention in organizational economics since
\citet{coase1937}. Williamson's Transaction Cost Economics
(TCE)~\citep{williamson1985} identified asset specificity, uncertainty,
and frequency as the conditions under which hierarchical governance
dominates market exchange. The property-rights theory of the
firm~\citep{hart1995} located the firm boundary at the allocation of
residual control rights over non-contractible decisions. Across these
traditions, a common premise is that market contracting between
strangers is expensive: search is costly, negotiation is costly, and
enforcement---through courts, arbitrators, or platforms---is costly,
slow, and jurisdiction-dependent.

This paper asks what happens to that boundary when one of the
underlying costs---contract enforcement between mutually distrusting
parties---collapses to a fixed, jurisdiction-independent capital
lockup. The collapse is not a thought experiment: the companion
mechanism paper presents the mechanism with
formal proofs of cooperation as the unique surviving Nash
equilibrium, and the companion implementation paper
presents a verified Solidity implementation audited through
TLA\textsuperscript{+} model checking, property-based fuzzing,
symbolic execution, and specification checking.

We treat that primitive as exogenous. The contribution of this paper
is institutional-economic: given the primitive, what follows for the
theory of the firm? We argue three things. First, the Coasean
threshold shifts inward---narrowly, in a precise sense we make
explicit. Second, the firm's various functions dissolve at different
rates: coordination dissolves where the primitive is applicable;
capital aggregation, organizational knowledge, risk pooling, and legal
personhood persist. Third, the practical difficulty of communicating
this shift is itself diagnostic: the institutions whose necessity
becomes contingent are precisely the institutions whose necessity has
achieved hegemonic invisibility.

\paragraph{Paper organization.}
\Cref{sec:primitive} summarizes the settlement primitive, citing the
companion paper for the formal development.
\Cref{sec:threshold} develops the Coasean threshold shift.
\Cref{sec:firm-types} distinguishes the coordination firm from the
capital firm and locates the dissolution.
\Cref{sec:token} examines token denomination as a coordination signal
that replaces standing platform credentials.
\Cref{sec:hegemony} offers a Gramscian reading of why the shift is
hard to see.
\Cref{sec:conclusion} concludes.

% ════════════════════════════════════════════════════════════════════
\section{The Settlement Primitive}
\label{sec:primitive}

We summarize here only the properties of the primitive that the
present argument uses. The full definitions and theorems are
in the mechanism paper; the implementation and verification
in the implementation paper.

\paragraph{Asymmetric bonding.}
For a transaction with payment $\pay > 0$ and cumulative upstream
value $\cumval \geq \pay$, the buyer locks $2\pay$ and the seller
locks $2\cumval$. Only the buyer can trigger resolution. On
resolution, the buyer recovers $\pay$ (net cost: $-\pay$) and the
seller recovers $2\cumval + \pay$ (net gain: $+\pay$). On
non-resolution, both bonds remain locked indefinitely.

\paragraph{Equilibrium properties.}
In the resulting one-shot bonding game, cooperation weakly dominates
defection for all parties, and the cooperative profile is the unique
outcome surviving iterated elimination of weakly dominated
strategies. The result extends to $N$-party process trees through
\emph{progressive collateralization}, with the seller's
risk-to-reward ratio growing with chain depth.

\paragraph{Atomic resolution and peer enforcement.}
All orders within a process resolve atomically---all or nothing---which
induces a weakest-link subgame among sellers. This produces an
endogenous coordination pressure analogous to peer enforcement in
joint-liability microfinance, but obtained without repeated
interaction, local information, or social punishment technology.

\paragraph{No escape hatches.}
The primitive has no timeout, no admin override, no governance
mechanism, no protocol fee, and no upgrade path. The companion paper
proves that any such escape hatch weakens the equilibrium.

\paragraph{What the primitive does \emph{not} provide.}
For the institutional argument that follows, two non-properties matter
as much as the properties. The primitive does not provide capital
aggregation across projects, does not provide pooled risk-bearing
across uncorrelated ventures, and does not provide legal personhood
for off-chain liability. We return to these absences in
\Cref{sec:firm-types}.

% ════════════════════════════════════════════════════════════════════
\section{The Coasean Threshold Shift}
\label{sec:threshold}

Coase's central insight is that firms exist because market
transaction costs---searching for partners, negotiating terms,
enforcing agreements---sometimes exceed the cost of hierarchical
management. Williamson refined this into TCE, identifying three
conditions under which hierarchical governance dominates market
exchange: high asset specificity, high uncertainty, and high
frequency.

The settlement primitive changes the calculus by collapsing one of
these costs to a fixed, predictable capital lockup. Consider the
three transaction costs separately.

\paragraph{Search costs.}
Public on-chain signals (settlement history, accepted-token lists,
attested capabilities) allow autonomous discovery through shared
coordination graphs. For agents able to read blockchain state,
search cost approaches zero. The substantive change is not that
search is cheaper---platforms already make search cheap---but that
the search index is non-proprietary: no platform owns the discovery
function, and the data is not extracted as rent.

\paragraph{Negotiation costs.}
Off-chain commitment building (EIP-712 typed structured data) lets
parties iterate on terms without gas costs. The negotiation is
bilateral and direct; counterparties can discover and iterate through
shared graph visibility rather than through a standing intermediary's
matching function. The substantive change here is asymmetric: the
\emph{cost} of negotiating is roughly unchanged from web2 baselines,
but the \emph{counterparty space} from which one can negotiate is no
longer mediated by a platform's terms of service.

\paragraph{Enforcement costs.}
This is the decisive variable. In the pre-primitive world, enforcement
requires courts (slow, expensive, jurisdiction-dependent), arbitrators
(trusted third parties), or standing intermediaries that internalize
coordination and charge for it. The primitive prices enforcement at
exactly $2\times$ the transaction value---a fixed, predictable,
jurisdiction-independent cost requiring no third party.

\begin{proposition}[Coasean Threshold Shift]
\label{prop:coasean-shift}
For any transaction where $2\pay$ (the capital lockup cost) is less
than the firm's coordination overhead for the same transaction, the
rational organizational choice is market exchange via the primitive
rather than hierarchical integration.
\end{proposition}

\Cref{prop:coasean-shift} is a comparative-statics statement, not a
sweeping institutional prediction. It identifies a class of
transactions for which the make-or-buy boundary tilts toward
\emph{buy}---specifically, transactions where the time-weighted
opportunity cost of locked capital is lower than the standing
overhead of internalizing the coordination function. The class is
non-empty (most retail-scale transactions in mature financial
environments satisfy it) and bounded (very-high-frequency, very-low-margin,
or very-long-duration transactions do not). The proposition does not
claim that all transactions migrate; it claims that the migration
threshold is now defined by capital opportunity cost rather than by
the cost of standing institutional infrastructure.

\paragraph{Engagement with TCE.}
Williamson's three drivers of vertical integration deserve direct
treatment. \emph{Asset specificity} continues to favor hierarchy: a
specialized asset whose value is contingent on a specific
counterparty cannot be efficiently coordinated through a sequence of
short-lived bonded transactions, since the holdup risk is in the
asset's pre-transaction sunk investment, not in the transaction
itself. \emph{Uncertainty} cuts both ways: contract incompleteness
favors hierarchy, but the primitive's evidence trail
(timestamped commitments, signed disclosures) reduces some classes of
uncertainty---specifically those that hierarchical authority
historically resolved by managerial fiat. \emph{Frequency} is the
condition under which the shift is most pronounced: high-frequency
exchanges with low asset specificity are exactly where the standing
firm's overhead is hardest to amortize and the primitive's per-transaction
capital lockup is easiest to absorb.

The proposition therefore narrows rather than overturns TCE. It
identifies a region of the asset-specificity / uncertainty / frequency
space where hierarchical governance was previously the equilibrium
outcome but is no longer.

% ════════════════════════════════════════════════════════════════════
\section{From Coordination Firm to Capital Firm}
\label{sec:firm-types}

The implication of \Cref{prop:coasean-shift} is structural: the
boundary of the firm shifts inward in a precise region. The stronger
implication is organizational rather than merely contractual. When
enforcement is priced directly in bonded capital, the standing firm
ceases to be the default container of production for the class of
transactions identified above. What emerges in that region is a
\emph{transaction-scoped institution}: a temporary assembly of
directly bonded counterparties that forms around a process,
coordinates for the duration of that process, and dissolves at
settlement.

The dissolution is narrower than \Cref{prop:coasean-shift} alone may
suggest. The firm historically bundles at least five distinct
functions:
\begin{enumerate}[label=(\roman*)]
  \item \textbf{Coordination of production} --- aligning the activities
    of multiple value-adders to produce a finished output.
  \item \textbf{Capital aggregation} --- pooling capital across
    projects to fund investment that no single participant can
    finance alone.
  \item \textbf{Organizational knowledge} --- the accumulated tacit
    procedural and relational knowledge that is hard to codify.
  \item \textbf{Risk pooling} --- diversifying across uncorrelated
    ventures to smooth participant income.
  \item \textbf{Legal personhood} --- providing a contractable entity
    for liability, regulatory compliance, and external counterparty
    relationships.
\end{enumerate}

The settlement primitive addresses (i) directly. It does not
address (ii)--(v). The remaining functions persist as requirements
that some durable institution must satisfy; they likely reconstitute
around token-denominated treasury communities whose members
repeatedly participate in compatible process trees---structures that
resemble firms at the meta-layer.

The precise claim is therefore: the \emph{coordination firm}, whose
value is internal management of a bounded labor pool, loses its
economic rationale where the primitive is applicable; the
\emph{capital firm}, whose value is durable asset ownership and
risk-bearing, persists in revised wrappers.

\paragraph{Key components and individually addressable capability.}
Each seller node in a process tree corresponds to what we term a
\emph{key component}---the irreducible asset or capability that
provides competitive advantage in a specific market. In the firm
model, key components are embedded inside organizational hierarchies:
the cook's skill, the kitchen's equipment, the courier's vehicle are
all subsumed under the firm's legal identity. In the primitive
model, each key component is represented directly by a seller wallet
address. The bonded process tree does not merely flatten the firm---it
makes each value-adding capability individually addressable,
individually compensated, and individually accountable. The
productive asset \emph{is} the identity.

Nothing in the primitive requires that a node be a natural person or
a corporation. Any actor capable of controlling an address and
bearing capital risk---a human, a software agent, a machine operator,
or a treasury---can occupy a node in the process tree. The
accountable unit is the bonded counterparty, not the firm.

\paragraph{What this means for the property-rights theory.}
The property-rights tradition~\citep{hart1995} locates the firm
boundary at the allocation of residual control rights over
non-contractible decisions. Under the primitive, certain previously
non-contractible decisions become contractible: schema-typed
attestations and atomic-resolution conditions encode behaviors that
historically required human discretion (acceptance, dispute,
quality verification). The class of residual non-contractible
decisions therefore shrinks. The property-rights logic still applies
to what remains; the boundary moves because the set of contractible
behaviors expands.

% ════════════════════════════════════════════════════════════════════
\section{Token Denomination as Coordination Signal}
\label{sec:token}

In the primitive, the choice of denomination token is itself a
coordination signal. The settlement layer is token-agnostic: any
ERC-20 can serve as the bond and payment currency. This means:

\begin{itemize}
  \item A stablecoin signals legal-system alignment.
  \item A DAO governance token signals community membership.
  \item A commodity-backed token signals value anchoring.
  \item Accepting a specific token \emph{is} the seller's identity
    declaration---an accepted-token list is a brand.
\end{itemize}

Settlement velocity in a given token replaces traditional reputation
metrics (star ratings, review scores). The on-chain record of resolved
orders denominated in token $\tau$ is a public, tamper-proof signal of
the seller's reliability within the $\tau$ economy. No platform needs
to aggregate, curate, or attest this signal; it is a natural
byproduct of protocol usage.

This matters institutionally because reputation infrastructure is one
of the durable rents that platforms extract. A standing platform's
credentialing function---certifying that a counterparty is reliable,
safe, or qualified---is one of the structural reasons platforms
persist even after their matching function is technically replicable.
When reputation is a public byproduct of settlement rather than a
proprietary platform asset, the platform's claim to that rent
weakens.

% ════════════════════════════════════════════════════════════════════
\section{Cultural Hegemony and the Visibility of Coordination}
\label{sec:hegemony}

\citet{gramsci1971} observed that the most durable forms of power are
those embedded so deeply in common sense that they become invisible.
He called this \emph{cultural hegemony}: the dominant arrangements
rule not only through coercion but by making themselves appear
natural, inevitable, and beyond question. The concept applies
directly to contemporary coordination infrastructure.

The platform economy illustrates Gramscian hegemony precisely. The
proposition that two strangers need a platform to transact safely
feels like common sense. It is not. It is a contingent arrangement
that arose because no economic primitive existed for safe direct
exchange between strangers. In Gramscian terms, the platform has
achieved hegemonic status: its necessity is no longer argued for---it
is assumed.

Similarly, the proposition that production requires firms (the
Coasean assumption examined in \Cref{prop:coasean-shift}) and that
cross-border trade requires corridors controlled by state or alliance
actors (ports, railways, digital infrastructure) have both achieved
the same hegemonic status. These are presented as structural
necessities rather than contingent design choices.

The primitive makes the coordination function \emph{visible} by
separating it from the entities that currently perform it. The
platform's coordination function (matching, enforcement, dispute
resolution) becomes a public smart contract. The firm's coordination
function (search, negotiation, enforcement) becomes three reducible
costs. The corridor's coordination function (settlement, standards,
access control) becomes a bilateral mechanism. In each case, the
entity does not need to be \emph{replaced}---it needs to be
\emph{seen} as a design choice rather than a fact of nature. Once the
coordination function is visible and separable, the architecture that
makes the entity structurally unnecessary follows.

This analysis explains a practical difficulty in communicating the
implications of the primitive. Proposing to remove structures that
most people do not know are there is qualitatively different from
proposing a better version of the same structure. The challenge is
not engineering. It is making existing infrastructure legible as
architecture rather than gravity.

\paragraph{A scope caveat.}
The Gramscian framing identifies a barrier to recognition; it does
not, by itself, predict which institutions will dissolve and which
will persist. The threshold proposition (\Cref{prop:coasean-shift})
and the firm-type distinction (\Cref{sec:firm-types}) do that work.
The hegemony argument explains why the work is needed: not because
the institutional consequence is empirically subtle, but because the
incumbent arrangement is presupposed in the language used to reason
about alternatives.

% ════════════════════════════════════════════════════════════════════
\section{Conclusion}
\label{sec:conclusion}

The Coasean threshold is not a conjecture---it is observable.
Contemporary coordination stacks often extract 15--30\% of transaction
value as coordination rent (Uber~25\%, DoorDash~30\%, Airbnb~14\%,
Amazon Marketplace~15\%). Credit card networks layer 2--3\%
interchange plus double-digit annual interest on revolving balances.
Financial markets intermediate trillions in notional value against a
fraction of that in real economic exchange. Asymmetric bonding prices
the cost of trust at $2\times$ the transaction value in temporarily
locked (not spent) capital, with zero recurring extraction.

For any economic activity where the coordination rent exceeds the
opportunity cost of locked capital, the standing firm---as Coase
conceived it---is no longer the uniquely efficient unit of economic
organization in that activity. The stronger consequence is the
emergence of transaction-scoped institutions: temporary assemblies of
directly bonded contributors that form around a process and dissolve
at settlement. The settlement primitive does not abolish coordination;
it changes the unit of coordination from the standing firm to the
bonded process---in the regions where it is applicable.

The argument leaves several questions open. First, the empirical
question: which industries and which transaction classes lie above
the threshold today, and at what rate does the threshold shift as
on-chain settlement matures? Second, the institutional question: how
do the firm functions that the primitive does not address---capital
aggregation, organizational knowledge, risk pooling, legal
personhood---reconstitute around the primitive without re-creating
the coordination firm by the back door? Third, the political question:
who pays the transition costs of replacing standing intermediaries
with bilateral mechanisms, and what coalitions form to slow or
accelerate the shift? Each is empirical or institutional rather than
mechanism-theoretic; the present paper has limited itself to
identifying that the questions are now well-posed.

\section*{Acknowledgments}
I thank the readers and collaborators who pressed me to separate the
mechanism-design contribution (developed in the companion paper) from
the institutional argument made here. The two were entangled in
earlier drafts in a way that did neither justice.


% ════════════════════════════════════════════════════════════════════
% References
% ════════════════════════════════════════════════════════════════════

\begin{thebibliography}{99}

\bibitem[Coase(1937)]{coase1937}
Coase, R.~H.
\newblock The Nature of the Firm.
\newblock \emph{Economica}, 4(16):386--405, 1937.

\bibitem[Gramsci(1971)]{gramsci1971}
Gramsci, A.
\newblock \emph{Selections from the Prison Notebooks}.
\newblock International Publishers, New York, 1971.
\newblock Edited and translated by Q.~Hoare and G.~Nowell-Smith.

\bibitem[Hart(1995)]{hart1995}
Hart, O.
\newblock \emph{Firms, Contracts, and Financial Structure}.
\newblock Oxford University Press, 1995.

\bibitem[Williamson(1985)]{williamson1985}
Williamson, O.~E.
\newblock \emph{The Economic Institutions of Capitalism}.
\newblock Free Press, New York, 1985.

\end{thebibliography}

\end{document}
